% Part: computability
% Chapter: computability-theory
% Section: non-comp-set

\documentclass[../../../include/open-logic-section]{subfiles}

\begin{document}

\olfileid{cmp}{thy}{ncp}
\olsection{توجد مجموعات غير قابلة للحساب}

رأينا أعلاه أن كل مجموعة قابلة للحساب قابلة للتعداد حسابيًا. فهل يصح
العكس؟ يبين الآتي أنه لا يصح بوجه عام.

\begin{thm}\ollabel{thm:K-0}
لتكن $K_0$ المجموعة $\Setabs{\tuple{e, x}}{\cfind{e}(x) \fdefined}$.
عندئذ تكون $K_0$ قابلة للتعداد حسابيًا، ولكنها غير قابلة للحساب.
\end{thm}

\begin{proof}
لرؤية أن $K_0$ قابلة للتعداد حسابيًا، لاحظ أنها مجال الدالة~$f$
المعرفة بـ
\[
f(z) = \umin{y}{(\len{z} = 2 \land T((z)_0, (z)_1, y))}.
\]
فإذا كانت $\cfind{e}(x)$ معرفة، وجدت $f(\tuple{e,x})$ متتالية حساب
متوقف؛ وإذا كانت $\cfind{e}(x)$ غير معرفة، كانت
$f(\tuple{e,x})$ كذلك؛ وإذا لم يكن $z$ يرمز زوجًا أصلًا، كانت $f(z)$
غير معرفة أيضًا.

أما كون $K_0$ غير قابلة للحساب، فهو بالضبط عدم قابلية مشكلة التوقف
للقرار، \olref[hlt]{thm:halting-problem}.
\end{proof}

المجموعة $K_0$ هي مجموعة الأزواج $\tuple{e,x}$ التي تكون عندها
$\cfind{e}(x)\fdefined$؛ أي إن $\tuple{e,x}\in K_0$ إذا وفقط إذا كانت
$\cfind{e}$ معرفة، أو يتوقف حسابها، عند الدخل~$x$، ولذلك تسمى أيضًا
«مجموعة التوقف». أما المجموعة
$K=\Setabs{e}{\cfind{e}(e)\fdefined}$ فهي «مجموعة التوقف الذاتي»،
وكثيرًا ما تستعمل مثالًا معياريًا لمجموعة غير قابلة للقرار.

\begin{thm}\ollabel{thm:K}
مجموعة التوقف الذاتي $K=\Setabs{e}{\cfind{e}(e)\fdefined}$
!!{c.e.} ولكنها غير قابلة للقرار.
\end{thm}

\begin{proof}
  افترض أن $K$ قابلة للقرار، أي إن دالتها المميزة $\Char{K}$ قابلة
  للحساب. لتكن
  \[d(e) = \begin{cases}
  1 & \text{إذا كانت\/ $\Char{K}(e) = 0$}\\
  \fundefined & \text{في غير ذلك.}
  \end{cases}
  \]
  وليكن $k$ فهرس~$d$، أي $d\simeq\cfind{k}$. عندئذ
  $d(k)\simeq\cfind{k}(k)$. وهذا يناقض حقيقة أن
  $d(k)\fdefined$ إذا وفقط إذا كانت $\cfind{k}(k)\fundefined$، وهي
  حقيقة تتبع من تعريف~$d$.

  و$K$ مجال $f(x)=\umin{y}{T(x,x,y)}$، ولذلك هي !!{c.e.}.
\end{proof}

\end{document}
